\label{sec:4}

Compelling argument demonstrating success criteria met, or well-justified explanation of different direction taken.
Excellent critical thought and interpretation which substantiate any claims of success.

\section{Evaluation Setup}
\label{sec:4.1}

% TODO


\section{System Performance}
\label{sec:4.2}

Reinforcement Learning with multiple agents requires millions of datapoints for updating the polcies and an equal number to evaluate the language and policy.
This section quantitivley measures the performance gains through the custom systems developed.
Specifically, it analyses the steps per second of the rust simulation developed and the memory efficiency of the \lstinline[language=Python]|RolloutBuffer|.
Benchmarks were evaludated on NVIDIA RTX 5060 Ti and AMD Ryzen 7 3700X.

\subsection{Vectorised Environment Throughput}
\label{sec:4.2.1}

A major issue with Python running multiple threads concurrently is the GIL.
While Python's builtin library \lstinline[language=Python]|multiprocessing| can, in theory, bypass this limitation, it introduces considerable overhead.
For this RL project, all threads require synchronisation at every simulation timestep.

To justify developing the simulation in Rust rather than Python, a simple mock simulation was developed in Python.
This mock simulation was compared to the vectorised Rust backend.
Environment steps per second (SPS) was calculated to analyse the throughput of each of the simulations.
Number of worlds parallised was scaled ($W \in \{2^0, 2^1, ..., 2^8\}$) to measure concurrency efficiency.

\begin{figure}[htbp]
    \centering
    \includegraphics[width=0.95\textwidth]{sim-comparison.png}
    \caption{Concurrent Environment Throughput Performance. Comparison of Python multiprocessing and Rust with Rayon across an exponentially increasing number of parallel worlds}
    \label{fig:sim-time-comparison}
\end{figure}
% TODO: Test for how long?

As shown in figure \ref{fig:sim-time-comparison}, the Python baseline completely fails to match the Rust simulation.
Despite being able to bypass the GIL and using multiple CPU cores, the serialisation and the Inter-Process Communication overhead required to transfer the necassary data between isolated processes at every time step heavily bottleneck the throughput.

In contrast, the Rust is able utilise \lstinline[language=Rust]|Rayon|'s thread-pool to efficiently manage the parallel world steps nativly in C-space.
After a minor scheduling overhead at $W=2$, throughput scales linearly with number of worlds parallised in $\log_2$.
By using Rust's native Rayon threading, the sequential fraction of the program is minimised per Amdahl's Law allowing near-perfect scaling that saturates the CPU cores.

\subsection{Rollout Buffer Efficiency}

A secondary major potential bottleneck is transferring the results of the rollout phase to the training phase (across CPU-GPU boundary).
To evaluate the \lstinline[language=Python]|RolloutBuffer|, the memory allocation and transfer overhead is compared to a baseline method (appending to Python lists and executing \lstinline[language=Python]|torch.stack| after the rollout phase).

\begin{table}[h]
    \centering
    \begin{tabular}{|l|c|c|c|}
    \hline
    \textbf{Memory Strategy} & \textbf{Allocation Time (s)} & \textbf{Peak VRAM Overhead (MB)} \\ \hline
    \textbf{Dynamic Python Lists} & $6.47 \pm 0.18$ & $1213.45$ \\ \hline
    \textbf{Pre-allocated RolloutBuffer} & $6.79 \pm 0.23$ & $598.23$ \\ \hline
    \end{tabular}
    \caption{Memory Management Overhead per Epoch: $T=20,000$, $W=32$}
    \label{tab:rolloutbuffer-eval}
\end{table}

As shown in \ref{tab:rolloutbuffer-eval}, the \lstinline[language=Python]|RolloutBuffer| method reduces VRAM by over $50\%$.
This is as Python's dynamic arrays often require copying their contents from one memory location to another to increase their size.
Whilst there is a $5\%$ improvement in allocation time when using Python's dynamic list, the large improvement in VRAM will enable scaling this method to millions of datapoints.


\section{Offline Pre-training}
\label{sec:4.3}

\subsection{Reconstruction and Commitment Quality}
\label{sec:4.3.1}

\subsection{Initial Codebook Health}
\label{sec:4.3.2}


\section{Online MARL Task Performance}
\label{sec:4.4}

\subsection{Baseline Comparisons}
\label{sec:4.4.1}

\subsection{Bandwidth-Performance Trade-off}
\label{sec:4.4.2}


\section{Analysis of Emergent Linguistics}
\label{sec:4.5}

\subsection{Spontaneous Semantic Compression}
\label{sec:4.5.1}

\subsection{Intent and Semantic Mapping}
\label{sec:4.5.2}

\subsection{Decoding the HQ-VAE}
\label{sec:4.5.3}


\section{Architectural Ablation Studies}
\label{sec:4.6}

\subsection{Reflection Strategies - Are They Needed?}
\label{sec:4.6.1}


\subsection{VQ-VAE vs SQ-VAE}
\label{sec:4.6.2}